\section{Remaining questions}

For repeated exponents $a\geq10$, the next question is whether the
symmetric cubic $f_{a,b}$ always has a noninteger root on the admissible
square-discriminant pairs with
\[
1\leq b\leq\left\lfloor\frac{2(a-4)(a-2)}{3a+4}\right\rfloor.
\]
Their smaller factor satisfies $d\geq3a+4$, with index $j\geq3$.
The two divisor criteria give finite problems for each fixed $a$ or each
fixed $j$, but the resulting family over arbitrary parameters is not settled.
The cases with three pairwise distinct exponents, and nonsquarefree
exponent vectors with more than three prime factors, require further
arguments. The results here therefore do not establish nonintegrality
for every integer with at least three distinct prime factors.

A useful direction is to combine the arithmetic conditions on $(d,e)$
with a sign or congruence obstruction for the symmetric cubic. The
distinction between the two blocks is essential: an integral
antisymmetric block occurs on the boundary family, while an integral
symmetric cubic occurs at $(1,1,2)$.
